%
%Map ID hợp nhất 4 mức độ.
%
%Nếu là ID5 sẽ theo định đạng: %[Tham số 1 Tham số 2 Tham số 3 Tham số 4 Tham số 5]. Ví dụ: %[1D2B3]
%Nếu là ID6 sẽ theo định đạng: %[Tham số 1 Tham số 2 Tham số 3 Tham số 4 Tham số 5 - Tham số 6]. Ví dụ: %[1D2B3-1]
%Có thể thay đổi nội dung của mức độ nhưng vị trí của mức độ trong ID không đổi (thông số thứ 3 từ trái qua).
%Cú pháp của thông số: [Giá trị] Mô tả thông số. Ví dụ: [0] Lớp 10 thì 0 là giá trị để lưu vào ID.
%
%Chú ý: Các dấu gạch ngang không được thay đổi.
%
%Cấu hình tên các tham số
%
Tên tham số 1: Lớp
Tên tham số 2: Môn
Tên tham số 3: Chương
Tên tham số 4: Mức độ
Tên tham số 5: Bài
Tên tham số 6: Dạng
%
%Cấu hình chi tiết ID
%
%%Cấu hình mức độ dùng chung.
[N] Nhận biết
[H] Thông hiểu
[V] Vận dụng
[C] Vận dụng cao
%
%Cấu hình nội dung
%
-[0] Lớp 10
----[D] Đại số/Thống kê 10
-------[1] Mệnh đề. Tập hợp
----------[1] Mệnh đề
-------------[1] Xác định mệnh đề, mệnh đề chứa biến
-------------[2] Tính đúng-sai của mệnh đề
-------------[3] Phủ định của một mệnh đề
-------------[4] Mệnh đề kéo theo, mệnh đề đảo, mệnh đề tương đương
-------------[5] Mệnh đề với mọi, tồn tại (và phủ định chúng)
-------------[6] Áp dụng mệnh đề vào suy luận có lí
----------[2] Tập hợp
-------------[1] Tập hợp và phần tử của tập hợp
-------------[2] Tập hợp con. Hai tập hợp bằng nhau
-------------[3] Ký hiệu khoảng, đoạn, nửa khoảng
----------[3] Các phép toán tập hợp
-------------[1] Giao và hợp của hai tập hợp (rời rạc)
-------------[2] Hiệu và phần bù của hai tập hợp (rời rạc)
-------------[3] Giao và hợp (dùng đoạn, khoảng)
-------------[4] Hiệu và phần bù (dùng đoạn, khoảng)
-------------[5] Toán thực tế ứng dụng của tập hợp
-------[2] BPT và hệ BPT bậc nhất hai ẩn
----------[1] Bất phương trình bậc nhất hai ẩn
-------------[1] Các khái niệm về BPT bậc nhất hai ẩn
-------------[2] Miền nghiệm của BPT bậc nhất hai ẩn
-------------[3] Toán thực tế về BPT bậc nhất hai ẩn
----------[2] Hệ bất phương trình bậc nhất hai ẩn
-------------[1] Các khái niệm về Hệ BPT bậc nhất hai ẩn
-------------[2] Miền nghiệm của Hệ BPT bậc nhất hai ẩn
-------------[3] Toán thực tế về Hệ BPT bậc nhất hai ẩn
-------[3] Hàm số bậc hai và đồ thị
----------[1] Hàm số và đồ thị
-------------[1] Xác định một hàm số
-------------[2] Tập xác định của hàm số
-------------[3] Giá trị của hàm số
-------------[4] Đồ thị của hàm số
-------------[5] Tính đồng biến, nghịch biến
-------------[6] Tính chẵn, lẻ
-------------[7] Toán thực tế về hàm số
----------[2] Hàm số bậc hai
-------------[1] Xác định hàm số bậc hai
-------------[2] Bảng biến thiên, tính đơn điệu, max, min
-------------[3] Đồ thị của hàm số bậc hai
-------------[4] Bài toán về sự tương giao
-------------[5] Hàm số chứa dấu giá trị tuyệt đối
-------------[6] Toán thực tế ứng dụng hàm số bậc hai
-------[6] Thống kê
----------[1] Số gần đúng. Sai số
-------------[1] Tính và ước lượng sai số tuyệt đối, tương đối
-------------[2] Tính và xác định độ chính xác của kết quả
-------------[3] Quy tròn số gần đúng
-------------[4] Viết số gần đúng cho số đúng biết độ chính xác
----------[2] Mô tả và biểu diễn dữ liệu trên các bảng và biểu đồ
-------------[1] Đọc và phân tích thông tin trên bảng số liệu
-------------[2] Đọc và phân tích thông tin trên biểu đồ
-------------[3] Số liệu bất thường trên bảng số liệu
-------------[4] Số liệu bất thường trên biểu đồ
----------[3] Các số đặc trưng đo xu thế trung tâm của mẫu số liệu
-------------[1] Câu hỏi lý thuyết
-------------[2] Số trung bình cộng
-------------[3] Số trung vị
-------------[4] Tứ phân vị
-------------[5] Mốt
----------[4] Các số đặc trưng đo mức độ phân tán của mẫu số liệu
-------------[1] Câu hỏi lý thuyết
-------------[2] Khoảng biến thiên, khoảng tứ phân vị
-------------[3] Giá trị bất thường của mẫu số liệu
-------------[4] Phương sai, độ lệch chuẩn của mẫu số liệu
-------[7] Bất phương trình bậc 2 một ẩn
----------[1] Dấu của tam thức bậc 2
-------------[1] Xác định tam thức bậc 2
-------------[2] Dấu của tam thức bậc 2 và ứng dụng
-------------[3] Bài toán xét dấu biết BXD, đồ thị
-------------[4] Xét dấu biểu thức dạng tích, thương
-------------[5] Toán thực tế ứng dụng dấu tam thức bậc 2
----------[2] Giải bất phương trình bậc 2 một ẩn
-------------[1] Bất phương trình bậc 2 và ứng dụng
-------------[2] Giải bất phương trình bậc hai biết BXD, đồ thị
-------------[3] Bất phương trình dạng tích, thương
-------------[4] Hệ bất phương trình BPT bậc 2
-------------[5] Bất phương trình chứa căn, | · |
-------------[6] Bài toán có tham số m
-------------[7] Toán thực tế ứng dụng bất phương trình bậc 2 một ẩn
----------[3] Phương trình quy về phương trình bậc hai
-------------[1] Phương trình căn √︀f(x) = √︀g(x) và mở rộng
-------------[2] Phương trình căn √︀f(x) = g(x) và mở rộng
-------------[3] Phương trình căn thức có tham số
-------------[4] Phương trình căn thức (dạng khác)
-------------[5] Phương trình khác quy về phương trình bậc 2
-------------[6] Toán hình, toán thực tế ứng dụng phương trình quy về bậc 2
-------[8] Đại số tổ hợp
----------[1] Quy tắc cộng-quy tắc nhân
-------------[1] Bài toán chỉ sử dụng quy tắc cộng
-------------[2] Bài toán chỉ sử dụng quy tắc nhân
-------------[3] Bài toán kết hợp quy tắc cộng và quy tắc nhân
-------------[4] Bài toán dùng quy tắc bù trừ
-------------[5] Bài toán đếm số tự nhiên
-------------[6] Sơ đồ hình cây
----------[2] Hoán vị. Chỉnh hợp. Tổ hợp
-------------[1] Lý thuyết tổng hợp về P,C,A
-------------[2] Bài toán có biểu thức P,C,A
-------------[3] Bài toán đếm số tự nhiên
-------------[4] Bài toán chọn người
-------------[5] Bài toán chọn đối tượng khác
-------------[6] Bài toán có yếu tố hình học
-------------[7] Bài toán xếp chỗ (không tròn, không lặp)
-------------[8] Hoán vị bàn tròn
-------------[9] Hoán vị lặp
----------[3] Nhị thức Newton (Kể cả chuyên đề)
-------------[1] Các câu hỏi lý thuyết tổng hợp
-------------[2] Khai triển một nhị thức Newton
-------------[3] Tìm hệ số, số hạng trong khai triển bằng tam giác Pascal
-------------[4] Tìm hệ số, số hạng trong khai triển
-------------[5] Tính tổng nhờ khai triển nhị thức Newton
-------------[6] Toán tổ hợp có dùng nhị thức Newton
-------[0] Xác suất
----------[1] Không gian mẫu và biến cố
-------------[1] Các câu hỏi lý thuyết tổng hợp
-------------[2] Mô tả không gian mẫu, biến cố
-------------[3] Đếm phần tử không gian mẫu, biến cố
----------[2] Xác suất của biến cố
-------------[1] Các câu hỏi lý thuyết tổng hợp
-------------[2] Liên quan xúc xắc, đồng tiền (PP liệt kê)
-------------[3] Liên quan việc sắp xếp chỗ
-------------[4] Liên quan việc chọn người
-------------[5] Liên quan việc chọn đối tượng khác
-------------[6] Liên quan hình học
-------------[7] Liên quan việc đếm số
-------------[8] Liên quan bàn tròn hoặc hoán vị lặp 
-------------[9] Liên quan vấn đề khác
----[H] Hình học 10
-------[4] Hệ thức lượng trong tam giác
----------[1] Giá trị lượng giác của góc (0 − 180∘)
-------------[1] Xét dấu của biểu thức lượng giác
-------------[2] Tính các giá trị lượng giác
-------------[3] Biến đổi, rút gọn biểu thức lượng giác
----------[2] Định lý sin và định lý côsin
-------------[1] Bài toán chỉ dùng định lý Sin, Côsin
-------------[2] Bài toán có dùng công thức diện tích
-------------[3] Biến đổi, rút gọn biểu thức
-------------[4] Nhận dạng tam giác
----------[3] Giải tam giác và ứng dụng thực tế
-------------[1] Giải tam giác
-------------[2] Các ứng dụng thực tế
-------[5] Véctơ (chưa xét tọa độ)
----------[1] Khái niệm véctơ
-------------[1] Xác định một véctơ
-------------[2] Xét phương và hướng của các véctơ
-------------[3] Hai véctơ bằng nhau
-------------[4] Hai véctơ đối nhau
-------------[5] Độ dài của một véctơ
-------------[6] Toán thực tế áp dụng véctơ
----------[2] Tổng và hiệu của hai véctơ
-------------[1] Tính toán, thu gọn hiệu các véctơ
-------------[2] Tính đúng-sai của 1 đẳng thức véctơ
-------------[3] Tìm điểm nhờ đẳng thức véctơ
-------------[4] Tính độ dài của véctơ tổng, hiệu
-------------[5] Cực trị hình học
-------------[6] Toán thực tế áp dụng tổng hiệu véctơ
----------[3] Tích của một số với véctơ
-------------[1] Xác định k.→−a và độ dài của nó
-------------[2] Biến đổi, thu gọn 1 đẳng thức véctơ
-------------[3] Tìm điểm nhờ đẳng thức véctơ
-------------[4] Sự cùng phương của 2 véctơ và ứng dụng
-------------[5] Phân tích 1 véctơ theo 2 véctơ không cùng phương
-------------[6] Tính độ dài của véctơ tổng, hiệu
-------------[7] Tập hợp điểm
-------------[8] Cực trị hình học
-------------[9] Toán thực tế áp dụng tích 1 số với véctơ
----------[4] Tích vô hướng (chưa xét tọa độ)
-------------[1] Tích vô hướng, góc giữa 2 véctơ
-------------[2] Tìm góc nhờ tích vô hướng
-------------[3] Đẳng thức về tích vô hướng hoặc độ dài
-------------[4] Điều kiện vuông góc
-------------[5] Các bài toán tìm điểm và tập hợp điểm
-------------[6] Cực trị và chứng minh bất đẳng thức
-------------[7] Toán thực tế áp dụng tích vô hướng
-------[9] Phương pháp toạ độ trong mặt phẳng (Oxy)
----------[1] Toạ độ của véctơ
-------------[1] Tọa độ điểm, độ dài đại số của véctơ trên 1 trục
-------------[2] Phép toán véctơ (tổng, hiệu, tích với số) trong Oxy
-------------[3] Tọa độ điểm và véctơ trên hệ trục Oxy
-------------[4] Sự cùng phương của 2 véctơ và ứng dụng
-------------[5] Phân tích một véctơ theo 2 véctơ không cùng phương
-------------[6] Toán thực tế dùng hệ toạ độ
----------[2] Tích vô hướng (theo tọa độ)
-------------[1] Tích vô hướng, góc giữa 2 véctơ
-------------[2] Tìm góc nhờ tích vô hướng
-------------[3] Đẳng thức về tích vô hướng hoặc độ dài
-------------[4] Điều kiện vuông góc
-------------[5] Các bài toán tìm điểm và tập hợp điểm
-------------[6] Cực trị và chứng minh bất đẳng thức
-------------[7] Toán thực tế, liên môn
----------[3] Đường thẳng trong mặt phẳng toạ độ
-------------[1] Điểm, véctơ, hệ số góc của đường thẳng
-------------[2] Phương trình đường thẳng
-------------[3] Vị trí tương đối giữa hai đường thẳng
-------------[4] Bài toán về góc giữa hai đường thẳng
-------------[5] Bài toán về khoảng cách
-------------[6] Bài toán tìm điểm
-------------[7] Bài toán dùng cho tam giác, tứ giác
-------------[8] Bài toán thực tế, PP tọa độ hóa
-------------[9] Bài toán có dùng PT chính tắc
----------[4] Đường tròn trong mặt phẳng toạ độ
-------------[1] Tìm tâm, bán kính và điều kiện là đường tròn
-------------[2] Phương trình đường tròn
-------------[3] Phương trình tiếp tuyến của đường tròn
-------------[4] Vị trí tương đối liên quan đường tròn
-------------[5] Toán tổng hợp đường thẳng và đường tròn
-------------[6] Bài toán dùng cho tam giác, tứ giác
-------------[7] Bài toán thực tế, PP tọa độ hóa
----------[5] Ba đường conic trong mặt phẳng toạ độ và Chuyên đề 3
-------------[1] Xác định các yếu tố của elip
-------------[2] Phương trình chính tắc của elip
-------------[3] Bài toán điểm trên elip
-------------[4] Xác định các yếu tố của hypebol
-------------[5] Phương trình chính tắc của hypebol
-------------[6] Bài toán điểm trên hypebol
-------------[7] Xác định các yếu tố của parabol
-------------[8] Phương trình chính tắc của parabol
-------------[9] Bài toán điểm trên parabol
-------------[0] Bài toán tổng hợp/thực tế, PP tọa độ hóa 3 đường conic
----[C] Chuyên đề 10
-------[1] Hệ phương trình bậc nhất 3 ẩn và ứng dụng
----------[1] Hệ phương trình bậc nhất 3 ẩn và ứng dụng
-------------[1] Các khái niệm về Hệ PT bậc nhất 3 ẩn
-------------[2] Giải Hệ PT bậc nhất 3 ẩn
-------------[3] Toán thực tế ứng dụng Hệ PT bậc nhất 3 ẩn
-------[2] Phương pháp quy nạp toán học
----------[1] Phương pháp quy nạp toán học
-------------[1] Quy nạp chứng minh các đẳng thức/công thức/chia hết
-------------[2] Quy nạp chứng minh các bất đẳng thức
-[1] Lớp 11
----[D] Đại số & Giải tích/Thống kê 11
-------[1] Hàm số lượng giác và phương trình lượng giác
----------[1] Góc lượng giác
-------------[1] Câu hỏi lý thuyết
-------------[2] Chuyển đổi đơn vị độ và radian
-------------[3] Số đo của một góc lượng giác
-------------[4] Độ dài của một cung tròn
-------------[5] Đường tròn lượng giác và ứng dụng
-------------[6] Toán thực tế áp dụng góc lượng giác
----------[2] Giá trị lượng giác của một góc lượng giác
-------------[1] Câu hỏi lý thuyết
-------------[2] Xét dấu giá trị lượng giác. Tính giá trị lượng giác của một góc
-------------[3] Biến đổi, rút gọn biểu thức lượng giác; chứng minh đẳng thức lượng giác
-------------[4] Các góc lượng giác có liên quan đặc biệt: bù nhau, phụ nhau, đối nhau, hơn kém nhau π
-------------[5] Toán thực tế áp dụng giá trị của một góc lượng giác
----------[3] Các công thức lượng giác
-------------[1] Câu hỏi lý thuyết
-------------[2] Áp dụng công thức cộng
-------------[3] Áp dụng công thức nhân đôi - hạ bậc
-------------[4] Áp dụng công thức biến đổi tích <-> tổng
-------------[5] Kết hợp nhiều công thức lượng giác
-------------[6] Nhận dạng tam giác
-------------[7] Toán thực tế áp dụng công thức lượng giác
----------[4] Hàm số lượng giác và đồ thị
-------------[1] Câu hỏi lý thuyết
-------------[2] Tìm tập xác định
-------------[3] Xét tính đơn điệu
-------------[4] Xét tính chẵn, lẻ
-------------[5] Xét tính tuần hoàn, tìm chu kỳ
-------------[6] Tìm tập giá trị và min, max
-------------[7] Bảng biến thiên và đồ thị
-------------[8] Toán thực tế áp dụng hàm số lượng giác
----------[5] Phương trình lượng giác cơ bản
-------------[1] Câu hỏi lý thuyết. Khái niệm phương trình tương đương
-------------[2] Điều kiện có nghiệm
-------------[3] Phương trình cơ bản dùng Radian
-------------[4] Phương trình cơ bản dùng Độ
-------------[5] Phương trình đưa về dạng cơ bản
-------------[6] Toán thực tế áp dụng phương trình lượng giác
----------[6] [Giảm] Phương trình lượng giác thường gặp
-------------[1] Phương trình bậc n theo một hàm số lượng giác
-------------[2] Phương trình đẳng cấp bậc n đối với sinx và cosx
-------------[3] Phương trình bậc nhất đối với sinx và cosx
-------------[4] Phương trình đối xứng, phản đối xứng
-------------[5] Phương trình lượng giác không mẫu mực
-------------[6] Phương trình lượng giác có chứa ẩn ở mẫu số
-------------[7] Phương trình lượng giác có chứa tham số
-------------[8] Toán thực tế áp dụng phương trình lượng giác thường gặp
-------[2] Dãy số. Cấp số cộng. Cấp số nhân
----------[1] Dãy số
-------------[1] Câu hỏi lý thuyết
-------------[2] Số hạng tổng quát, biểu diễn dãy số
-------------[3] Tìm số hạng cụ thể của dãy số
-------------[4] Dãy số tăng, dãy số giảm
-------------[5] Dãy số bị chặn
-------------[6] Toán thực tế áp dụng dãy số
----------[2] Cấp số cộng
-------------[1] Câu hỏi lý thuyết
-------------[2] Nhận diện cấp số cộng, công sai d
-------------[3] Số hạng tổng quát của cấp số cộng
-------------[4] Tìm số hạng cụ thể trong cấp số cộng
-------------[5] Điều kiện để dãy số là cấp số cộng
-------------[6] Tính tổng của cấp số cộng
-------------[7] Toán thực tế áp dụng cấp số cộng
----------[3] Cấp số nhân
-------------[1] Câu hỏi lý thuyết
-------------[2] Nhận diện cấp số nhân, công bội q
-------------[3] Số hạng tổng quát của cấp số nhân
-------------[4] Tìm số hạng cụ thể trong cấp số nhân
-------------[5] Điều kiện để dãy số là cấp số nhân
-------------[6] Tính tổng của cấp số nhân
-------------[7] Kết hợp cấp số nhân và cấp số cộng
-------------[8] Toán thực tế áp dụng cấp số nhân
-------[3] Giới hạn. Hàm số liên tục
----------[1] Giới hạn của dãy số
-------------[1] Câu hỏi lý thuyết
-------------[2] Phương pháp đặt thừa số chung (lim hữu hạn)
-------------[3] Phương pháp lượng liên hợp (lim hữu hạn)
-------------[4] Giới hạn vô cực
-------------[5] Cấp số nhân lùi vô hạn
-------------[6] Toán thực tế áp dụng giới hạn của dãy số
----------[2] Giới hạn của hàm số
-------------[1] Câu hỏi lý thuyết
-------------[2] Thay số trực tiếp
-------------[3] PP đặt thừa số chung, kết quả hữu hạn
-------------[4] PP đặt thừa số chung, kết quả vô cực
-------------[5] PP lượng liên hợp, kết quả hữu hạn
-------------[6] PP lượng liên hợp, kết quả vô cực
-------------[7] Giới hạn một bên
-------------[8] Toán thực tế áp dụng giới hạn của hàm số
----------[3] Hàm số liên tục
-------------[1] Câu hỏi lý thuyết
-------------[2] Tính liên tục thể hiện qua đồ thị
-------------[3] Hàm số liên tục tại một điểm
-------------[4] Hàm số liên tục trên khoảng, đoạn
-------------[5] Bài toán phương trình có nghiệm
-------------[6] Toán thực tế áp dụng hàm số liên tục
-------[5] Các số đặc trưng đo xu thế trung tâm cho mẫu số liệu ghép nhóm
----------[1] Số trung bình và mốt của mẫu số liệu ghép nhóm
-------------[1] Câu hỏi lý thuyết
-------------[2] Mẫu số liệu ghép nhóm
-------------[3] Số trung bình
-------------[4] Mốt
----------[2] Trung vị và tứ phân vị của mẫu số liệu ghép nhóm
-------------[1] Câu hỏi lý thuyết
-------------[2] Trung vị
-------------[3] Tứ phân vị
-------[6] Hàm số mũ và hàm số lôgarít
----------[1] Phép tính luỹ thừa
-------------[1] Tính giá trị của biểu thức chứa lũy thừa
-------------[2] Biến đổi, rút gọn biểu thức chứa lũy thừa
-------------[3] Điều kiện cho luỹ thừa, căn thức
-------------[4] So sánh các lũy thừa
----------[2] Phép tính lôgarít
-------------[1] Tính giá trị biểu thức chứa lôgarít
-------------[2] Biến đổi, biểu diễn biểu thức chứa lôgarít
-------------[3] Rút gọn, chứng minh biểu thức lôgarít
-------------[4] Số e và bài toán lãi kép
-------------[5] Toán thực tế áp dụng phép tính lôgarít
----------[3] Hàm số mũ. Hàm số lôgarít
-------------[1] Câu hỏi lý thuyết hàm số lũy thừa, mũ, lôgarít
-------------[2] Tập xác định của hàm số
-------------[3] Sự biến thiên và đồ thị của hàm số mũ, lôgarít
-------------[4] So sánh các luỹ thừa và lôgarít
-------------[5] Toán thực tế áp dụng hàm số mũ, lôgarít
----------[4] Phương trình, bất phương trình mũ và lôgarít
-------------[1] Điều kiện có nghiệm
-------------[2] Phương trình mũ, lôgarít cơ bản
-------------[3] Bất phương trình mũ, lôgarít cơ bản
-------------[4] Phương trình mũ, lôgarít đưa về cùng cơ số
-------------[5] Bất phương trình mũ, lôgarít đưa về cùng cơ số
-------------[6] Toán thực tế áp dụng phương trình mũ, lôgarít
----------[5] [Giảm] Các phương pháp giải được giảm tải
-------------[1] Phương pháp đặt ẩn phụ cho PT mũ, lôgarít
-------------[2] Phương pháp lôgarít hóa, mũ cho PT mũ, lôgarít
-------------[3] Phương pháp hàm số, đánh giá cho PT mũ, lôgarít
-------------[4] Hệ PT mũ, lôgarít
-------------[5] Toán thực tế áp dụng phương trình mũ, lôgarít
-------[7] Đạo hàm
----------[1] Đạo hàm
-------------[1] Tính đạo hàm bằng định nghĩa
-------------[2] Số gia hàm số, số gia biến số
-------------[3] Ý nghĩa Hình học của đạo hàm
-------------[4] Ý nghĩa Vật lý của đạo hàm
-------------[5] Toán thực tế khác áp dụng định nghĩa đạo hàm
----------[2] Các quy tắc đạo hàm
-------------[1] Tính đạo hàm
-------------[2] Đẳng thức có y và y′
-------------[3] Tiếp tuyến tại một điểm
-------------[4] Tiếp tuyến biết trước hệ số góc
-------------[5] Tiếp tuyến chưa biết tiếp điểm và hệ số góc
-------------[6] Giới hạn hàm số lượng giác, hàm số mũ, lôgarít
-------------[7] Dùng đạo hàm cho nhị thức Newton
-------------[8] Toán thực tế áp dụng quy tắc đạo hàm
----------[3] Đạo hàm cấp hai
-------------[1] Tính đạo hàm cấp hai
-------------[2] Đẳng thức có y và (y′, y′′)
-------------[3] Toán thực tế và Ý nghĩa Vật lý của đạo hàm cấp hai
-------[9] Xác suất
----------[1] Biến cố giao và quy tắc nhân xác suất
-------------[1] Câu hỏi lí thuyết
-------------[2] Xác định và đếm số phần tử biến cố giao
-------------[3] Công thức nhân xác suất cho 2 biến cố độc lập
-------------[4] Tính xác suất biến cố giao bằng sơ đồ hình cây
----------[2] Biến cố hợp và quy tắc cộng xác suất
-------------[1] Câu hỏi lí thuyết
-------------[2] Mô tả biến cố hợp
-------------[3] Tính và xác định số phần tử biến cố hợp
-------------[4] Quy tắc cộng cho hai biến cố xung khắc
-------------[5] Quy tắc cộng cho hai biến cố bất kì
-------------[6] Tính xác suất bằng sơ đồ hình cây
----[H] Hình học 11
-------[4] Đường thẳng, mặt phẳng. Quan hệ song song trong không gian
----------[1] Điểm, đường thẳng và mặt phẳng
-------------[1] Câu hỏi lý thuyết
-------------[2] Hình biểu diễn của một hình không gian
-------------[3] Tìm giao tuyến của hai mặt phẳng
-------------[4] Tìm giao điểm của đường thẳng và mặt phẳng
-------------[5] Xác định thiết diện
-------------[6] Ba điểm thẳng hàng, ba đường thẳng đồng quy
-------------[7] Toán thực tế áp dụng điểm, đường thẳng và mặt phẳng
----------[2] Hai đường thẳng song song
-------------[1] Câu hỏi lý thuyết
-------------[2] Hai đường thẳng song song
-------------[3] Tìm giao tuyến bằng cách kẻ song song
-------------[4] Tìm giao điểm của đường thẳng và mặt phẳng
-------------[5] Xác định thiết diện bằng cách kẻ song song
-------------[6] Ba điểm thẳng hàng
-------------[7] Bài toán quỹ tích và điểm cố định
-------------[8] Toán thực tế áp dụng hai đường thẳng song song
----------[3] Đường thẳng và mặt phẳng song song
-------------[1] Câu hỏi lý thuyết
-------------[2] Đường thẳng song song với mặt phẳng
-------------[3] Tìm giao tuyến bằng cách kẻ song song
-------------[4] Tìm giao điểm của đường thẳng và mặt phẳng
-------------[5] Xác định thiết diện bằng cách kẻ song song
-------------[6] Ba điểm thẳng hàng
-------------[7] Bài toán quỹ tích và điểm cố định
-------------[8] Toán thực tế áp dụng đường thẳng song song mặt phẳng
----------[4] Hai mặt phẳng song song
-------------[1] Câu hỏi lý thuyết
-------------[2] Hai mặt phẳng song song
-------------[3] Chứng minh đường thẳng song song mặt phẳng
-------------[4] Xác định mặt phẳng đi qua một điểm và song song với một mặt phẳng
-------------[5] Xác định mặt phẳng chứa đường thẳng (hoặc đi qua hai điểm) và song song với một mặt phẳng
-------------[6] Bài toán tổng hợp
-------------[7] Toán thực tế áp dụng hai mặt phẳng song song
----------[5] Hình lăng trụ và hình hộp (xiên)
-------------[1] Câu hỏi lý thuyết
-------------[2] Bài toán về hình lăng trụ (xiên)
-------------[3] Bài toán về hình hộp (xiên)
-------------[4] Toán thực tế áp dụng hình lăng trụ và hình hộp
----------[6] Phép chiếu song song
-------------[1] Câu hỏi lý thuyết
-------------[2] Hình biểu diễn của một hình không gian
-------------[3] Xác định yếu tố song song
-------------[4] Xác định phương chiếu
-------------[5] Tính tỉ số đoạn thẳng, diện tích qua phép chiếu
-------[8] Quan hệ vuông góc trong không gian
----------[1] Hai đường thẳng vuông góc
-------------[1] Câu hỏi lí thuyết
-------------[2] Xác định hai đường thẳng vuông góc
-------------[3] Tìm góc giữa hai đường thẳng
-------------[4] Toán thực tế áp dụng hai đường thẳng vuông góc
----------[2] Đường thẳng vuông góc với mặt phẳng
-------------[1] Câu hỏi lí thuyết
-------------[2] Xác định hoặc chứng minh đường thẳng và mặt phẳng vuông góc
-------------[3] Xác định hoặc chứng minh hai đường thẳng vuông góc
-------------[4] Dựng mặt phẳng, tìm thiết diện
-------------[5] Hình chiếu vuông góc của một hình trên mặt phẳng (tìm điểm, tìm đoạn thẳng, tính diện tích)
-------------[6] Toán thực tế áp dụng đường thẳng vuông góc mặt phẳng
----------[3] Phép chiếu vuông góc
-------------[1] Lý thuyết về phép chiếu vuông góc
-------------[2] Hình chiếu vuông góc của đa giác trên mặt phẳng
-------------[3] Các bài toán thực tế áp dụng phép chiếu vuông góc
----------[4] Hai mặt phẳng vuông góc
-------------[1] Câu hỏi lí thuyết
-------------[2] Xác định/chứng minh đường thẳng vuông góc mặt phẳng, mặt phẳng vuông góc
-------------[3] Xác định góc giữa hai mặt phẳng
-------------[4] Dựng mặt phẳng, thiết diện
-------------[5] Nhận dạng và tính toán liên quan các hình thông dụng
-------------[6] Bài toán cho trước góc giữa d và (P)
-------------[7] Toán thực tế áp dụng hai mặt phẳng vuông góc
----------[5] Khoảng cách
-------------[1] Câu hỏi lí thuyết
-------------[2] Khoảng cách giữa 2 điểm, từ 1 điểm đến 1 đường thẳng
-------------[3] Khoảng cách từ một điểm đến một mặt phẳng
-------------[4] Khoảng cách giữa hai đường thẳng chéo nhau
-------------[5] Đường vuông góc chung của hai đường thẳng chéo nhau
-------------[6] Toán thực tế áp dụng khoảng cách
----------[6] Góc giữa đường thẳng và mặt phẳng. Góc nhị diện
-------------[1] Góc giữa đường thẳng và mặt phẳng
-------------[2] Góc nhị diện, góc phẳng nhị diện
-------------[3] Góc giữa 2 mặt phẳng, biết trước góc (d,(P))
-------------[4] Khoảng cách giữa điểm, đường, biết trước góc (d,(P))
-------------[5] Khoảng cách giữa điểm - mặt phẳng, biết trước góc (d,(P))
-------------[6] Khoảng cách giữa 2 đường chéo nhau, biết trước góc (d,(P))
-------------[7] Toán thực tế về góc đường thẳng, mặt phẳng, góc nhị diện
----------[7] Hình lăng trụ đứng. Hình chóp đều. Thể tích
-------------[1] Câu hỏi lí thuyết, công thức
-------------[2] Thể tích khối chóp tam giác
-------------[3] Thể tích khối chóp tứ giác
-------------[4] Thể tích khối lăng trụ tam giác
-------------[5] Thể tích khối lăng trụ tứ giác
-------------[6] Thể tích khối chóp cụt và các khối khác
-------------[7] Tỉ số thể tích
-------------[8] Ứng dụng thể tích tính góc, khoảng cách
-------------[9] Toán thực tế hình lăng trụ đứng, chóp đều, thể tích
----[C] Chuyên đề 11
-------[1] Phép biến hình phẳng
----------[1] Phép biến hình, phép dời hình và hai hình bằng nhau
-------------[1] Câu hỏi lý thuyết
-------------[2] Bài toán xác định một phép đặt tương ứng có là phép dời hình hay không?
-------------[3] Xác định ảnh khi thực hiện phép dời hình
----------[2] Phép tịnh tiến
-------------[1] Câu hỏi lý thuyết
-------------[2] Tìm ảnh hoặc tạo ảnh khi thực hiện phép tịnh tiến
-------------[3] Ứng dụng phép tịnh tiến
----------[3] Phép đối xứng trục
-------------[1] Câu hỏi lý thuyết
-------------[2] Tìm ảnh hoặc tạo ảnh khi thực hiện phép đối xứng trục
-------------[3] Xác định trục đối xứng và số trục đối xứng của một hình
-------------[4] Ứng dụng phép đối xứng trục
----------[4] Phép đối xứng tâm
-------------[1] Câu hỏi lý thuyết
-------------[2] Tìm ảnh, tạo ảnh khi thực hiện phép đối xứng tâm
-------------[3] Xác định hình có tâm đối xứng
-------------[4] Ứng dụng phép đối xứng tâm
----------[5] Phép quay
-------------[1] Câu hỏi lý thuyết
-------------[2] Xác định vị trí ảnh của điểm, hình khi thực hiện phép quay cho trước
-------------[3] Tìm tọa độ ảnh của điểm, phương trình của một đường thẳng khi thực hiện phép quay
-------------[4] Ứng dụng phép quay
----------[6] Phép vị tự
-------------[1] Câu hỏi lý thuyết
-------------[2] Xác định ảnh, tạo ảnh khi thực hiện phép vị tự
-------------[3] Tìm tâm vị tự của hai đường tròn
-------------[4] Ứng dụng phép vị tự
----------[7] Phép đồng dạng
-------------[1] Câu hỏi lý thuyết
-------------[2] Xác định ảnh, tạo ảnh khi thực hiện phép đồng dạng
-------[2] Lý thuyết đồ thị
----------[1] Đồ thị
-------------[1] Câu hỏi về đỉnh, cạnh của đồ thị
-------------[2] Câu hỏi về bậc của đồ thị
-------------[3] Câu hỏi tổng hợp
----------[2] Đường đi Euler và Harmilton
-------------[1] Đường đi Euler
-------------[2] Đường đi Harmilton
-------------[3] Câu hỏi tổng hợp
----------[3] Bài toán tìm đường đi ngắn nhất
-------------[1] Bài toán tìm đường đi ngắn nhất
-------------[2] Tổng hợp
-------[3] Một số yếu tố vẽ kỹ thuật
----------[1] Hình biểu diễn của một hình, khối
-------------[1] Lý thuyết về phép chiếu và hình biểu diễn song song
-------------[2] Lý thuyết về phép chiếu vuông góc
-------------[3] Lý thuyết về phép chiếu trục đo
-------------[4] Tổng hợp
----------[2] Bản vẽ kỹ thuật
-------------[1] Lý thuyết cơ bản về bản vẽ kỹ thuật
-------------[2] Phương pháp biểu diễn bản vẽ kỹ thuật
-------------[3] Tổng hợp
-[2] Lớp 12
----[D] Giải tích
-------[1] Ứng dụng đạo hàm để khảo sát hàm số
----------[1] Sự đồng biến và nghịch biến của hàm số
-------------[1] Xét tính đơn điệu của hàm số cho bởi công thức
-------------[2] Xét tính đơn điệu dựa vào bảng biến thiên, đồ thị
-------------[3] Tìm tham số $ m $ để hàm số đơn điệu
-------------[4] Ứng dụng tính đơn điệu để chứng minh bất đẳng thức, giải phương trình, bất phương trình, hệ phương trình
-------------[5] Toán thực tế ứng dụng sự đồng biến nghịch biến
----------[2] Cực trị của hàm số
-------------[1] Tìm cực trị của hàm số cho bởi công thức
-------------[2] Tìm cực trị dựa vào BBT, đồ thị
-------------[3] Tìm $ m $ để hàm số đạt cực trị tại $ 1 $ điểm $ x_0 $ cho trước
-------------[4] Tìm $ m $ để hàm số, đồ thị hàm số bậc ba có cực trị thỏa mãn điều kiện
-------------[5] Tìm $ m $ để hàm số, đồ thị hàm số trùng phương có cực trị thỏa mãn điều kiện
-------------[6] Tìm $ m $ để hàm số, đồ thị hàm số các hàm số khác có cực trị thỏa mãn điều kiện
-------------[7] Toán thực tế ứng dụng cực trị của hàm số
----------[3] Giá trị lớn nhất và giá trị nhỏ nhất của hàm số
-------------[1] GTLN, GTNN trên đoạn $ \left[a;b\right] $
-------------[2] GTLN, GTNN trên khoảng
-------------[3] Sử dụng các đánh giá, bất đẳng thức cổ điển
-------------[4] Ứng dụng GTNN, GTLN trong bài toán phương trình, bất phương trình, hệ phương trình
-------------[5] GTLN, GTNN hàm nhiều biến
-------------[6] Toán thực tế ứng dụng GTLN, GTNN của hàm số
----------[4] Đường tiệm cận
-------------[1] Bài toán xác định các đường tiệm cận của hàm số (không chứa tham số) hoặc biết BBT, đồ thị
-------------[2] Bài toán xác định các đường tiệm cận của hàm số có chứa tham số
-------------[3] Bài toán liên quan đến đồ thị hàm số và các đường tiệm cận
-------------[4] Toán thực tế ứng dụng tiệm cận
----------[5] Khảo sát sự biến thiên và vẽ đồ thị hàm số
-------------[1] Nhận dạng đồ thị
-------------[2] Các phép biến đổi đồ thị
-------------[3] Biện luận số giao điểm dựa vào đồ thị, bảng biến thiên
-------------[4] Sự tương giao của hai đồ thị (liên quan đến tọa độ giao điểm)
-------------[5] Đồ thị của hàm đạo hàm
-------------[6] Phương trình tiếp tuyến của đồ thị hàm số
-------------[7] Điểm đặc biệt của đồ thị hàm số
-------------[8] Toán thực tế ứng dụng khảo sát hàm số
-------[3] Các số đặc trưng đo mức độ phân tán cho mẫu số liệu ghép nhóm
----------[1] Khoảng biến thiên, khoảng tứ phân vị của mẫu số liệu ghép nhóm
-------------[1] Công thức lý thuyết
-------------[2] Tìm khoảng biến thiên
-------------[3] Tìm khoảng tứ phân vị
-------------[4] Câu hỏi tổng hợp
----------[2] Phương sai, độ lệch chuẩn của mẫu số liệu ghép nhóm
-------------[1] Công thức lý thuyết
-------------[2] Tìm phương sai, độ lệch chuẩn
-------------[3] Câu hỏi tổng hợp
-------[4] Nguyên hàm, tích phân và ứng dụng
----------[1] Nguyên hàm
-------------[1] Công thức lý thuyết
-------------[2] Nguyên hàm cơ bản đa thức, phân thức
-------------[3] Nguyên hàm cơ bản hàm lượng giác
-------------[4] Nguyên hàm cơ bản hàm mũ, luỹ thừa
-------------[5] Phương pháp đổi biến số cơ bản
-------------[6] Toán thực tế áp dụng nguyên hàm
----------[2] Tích phân
-------------[1] Công thức lý thuyết
-------------[2] Tích phân cơ bản đa thức, phân thức
-------------[3] Tích phân cơ bản hàm lượng giác
-------------[4] Tích phân cơ bản hàm mũ, luỹ thừa
-------------[5] Phương pháp đổi biến số cơ bản
-------------[6] Toán thực tế áp dụng nguyên hàm
----------[3] Ứng dụng thực tế và hình học của tích phân
-------------[1] Diện tích hình phẳng được giới hạn bởi các đồ thị
-------------[2] Bài toán thực tế sử dụng diện tích hình phẳng
-------------[3] Thể tích giới hạn bởi các đồ thị (tròn xoay)
-------------[4] Thể tích tính theo mặt cắt $S(x)$
-------------[5] Bài toán thực tế và ứng dụng thể tích tròn xoay, $S(x)$
-------[6] Một số yếu tố xác suất
----------[1] Xác suất có điều kiện
-------------[1] Công thức lý thuyết
-------------[2] Tính xác suất có điều kiện bằng công thức
-------------[3] Tính xác suất có điều kiện bằng sơ đồ cây
-------------[4] Bài toán tổng hợp
----------[2] Công thức xác suất toàn phần. Công thức Bayes
-------------[1] Công thức lý thuyết
-------------[2] Tính xác suất bằng công thức xác suất toàn phần
-------------[3] Tính xác suất bằng công thức xác suất Bayes
-------------[4] Bài toán tổng hợp
----[H] Hình học
-------[2] Tọa độ của véc-tơ trong không gian
----------[1] Véc-tơ và các phép toán véc-tơ trong không gian (chưa toạ độ hoá)
-------------[1] Công thức lý thuyết
-------------[2] Tổng, hiệu, tích một số với véc-tơ
-------------[3] Tích vô hướng và ứng dụng
-------------[4] Toán thực tế áp dụng các phép toán véc-tơ
----------[2] Toạ độ của véc-tơ và các công thức
-------------[1] Công thức lý thuyết
-------------[2] Tìm tọa độ điểm
-------------[3] Tìm tọa độ véc-tơ
-------------[4] Công thức toạ độ của tích vô hướng và ứng dụng
-------------[5] Công thức toạ độ của tích có hướng và ứng dụng
-------------[6] Toán thực tế áp dụng các phép toán toạ độ hoá véc-tơ}
-------[5] Phương trình mặt phẳng, đường thẳng, mặt cầu trong không gian $Oxyz$ 
----------[1] Phương trình mặt phẳng
-------------[1] Câu hỏi lý thuyết
-------------[2] Xác định véc-tơ pháp tuyến, cặp véc-tơ chỉ phương
-------------[3] Viết phương trình tổng quát mặt phẳng
-------------[4] Vị trí tương đối giữa hai mặt phẳng (song song, vuông góc)
-------------[5] Khoảng cách điểm tới mặt phẳng
-------------[6] Góc giữa hai mặt phẳng
-------------[7] Toán thực tế áp dụng phương trình mặt phẳng
----------[2] Phương trình đường thẳng trong không gian
-------------[1] Câu hỏi lý thuyết
-------------[2] Xác định véc-tơ chỉ phương, cặp véc-tơ pháp tuyến
-------------[3] Viết phương trình tổng quát, chính tắc, tham số đường thẳng
-------------[4] Vị trí tương đối giữa hai đường thẳng
-------------[5] Vị trí tương đối giữa đường thẳng và mặt phẳng
-------------[6] Khoảng cách điểm tới đường thẳng
-------------[7] Góc giữa hai đường thẳng, đường thẳng và mặt phẳng
-------------[8] Toán thực tế áp dụng phương trình đường thẳng
----------[3] Phương trình mặt cầu trong không gian
-------------[1] Câu hỏi lý thuyết
-------------[2] Xác định tâm, bán kính, đường kính mặt cầu
-------------[3] Viết phương trình tổng quát mặt cầu
-------------[4] Toán thực tế áp dụng phương trình mặt cầu
----[C] Chuyên đề
-------[1] Khối đa diện
----------[1] Khái niệm về khối đa diện
-------------[1] Nhận diện hình đa diện, khối đa diện
-------------[2] Xác định số đỉnh, cạnh, mặt bên của một khối đa diện
-------------[3] Phân chia, lắp ghép các khối đa diện
-------------[4] Phép biến hình trong không gian
----------[2] Khối đa diện lồi và khối đa diện đều
-------------[1] Nhận diện đa diện lồi
-------------[2] Nhận diện loại đa diện đều
-------------[3] Tính chất đối xứng
----------[3] Khái niệm về thể tích của khối đa diện
-------------[1] Diện tích xung quanh, diện tích toàn phần của khối đa diện
-------------[2] Tính thể tích các khối đa diện
-------------[3] Tỉ số thể tích
-------------[4] Các bài toán khác(góc, khoảng cách,...) liên quan đến thể tích khối đa diện
-------------[5] Bài toán thực tế về khối đa diện
-------------[6] Bài toán cực trị
-------[2] Mặt nón, mặt trụ, mặt cầu
----------[1] Khái niệm về mặt tròn xoay
-------------[1] Thể tích khối nón, khối trụ
-------------[2] Diện tích xung quanh, diện tích toàn phần, độ dài đường sinh, chiều cao, bán kính đáy, thiết diện
-------------[3] Khối tròn xoay nội tiếp, ngoại tiếp khối đa diện
-------------[4] Bài toán thực tế về khối nón, khối trụ
-------------[5] Bài toán cực trị về khối nón, khối trụ
----------[2] Mặt cầu
-------------[1] Bài toán sử dụng định nghĩa, tính chất, vị trí tương đối
-------------[2] Khối cầu ngoại tiếp khối đa diện
-------------[3] Khối cầu nội tiếp khối đa diện
-------------[4] Bài toán thực tế về khối cầu
-------------[5] Bài toán cực trị về khối cầu
-------------[6] Bài toán tổng hợp về khối nón, khối trụ, khối cầu
-------[3] Phương pháp tọa độ trong không gian
----------[1] Hệ tọa độ trong không gian
-------------[1] Tìm tọa độ điểm, véc-tơ liên quan đến hệ trục $Oxyz$
-------------[2] Tích vô hướng và ứng dụng
-------------[3] Phương trình mặt cầu (xác định tâm, bán kính, viết PT mặt cầu đơn giản, vị trí tương đối hai mặt cầu, điểm đến mặt cầu, đơn giản)
-------------[4] Các bài toán cực trị
----------[2] Phương trình mặt phẳng
-------------[1] Tích có hướng và ứng dụng
-------------[2] Xác định VTPT
-------------[3] Viết phương trình mặt phẳng
-------------[4] Tìm tọa độ điểm liên quan đến mặt phẳng
-------------[5] Góc
-------------[6] Khoảng cách
-------------[7] Vị trí tương đối giữa hai mặt phẳng, giữa mặt cầu và mặt phẳng
-------------[8] Các bài toán cực trị
----------[3] Phương trình đường thẳng trong không gian
-------------[1] Xác định VTCP
-------------[2] Viết phương trình đường thẳng
-------------[3] Tìm tọa độ điểm liên quan đến đường thẳng
-------------[4] Góc
-------------[5] Khoảng cách
-------------[6] Vị trí tương đối giữa hai đường thẳng, giữa đường thẳng và mặt phẳng
-------------[7] Bài toán liên quan giữa đường thẳng - mặt phẳng - mặt cầu
-------------[8] Các bài toán cực trị
----------[4] Ứng dụng của phương pháp tọa độ
-------------[1] Bài toán HHKG
-------------[2] Bài toán đại số
-[6] Lớp 6
----[C] Cánh Diều
-------[1] Số tự nhiên
----------[1] Tập hợp
----------[2] Tập hợp các số tự nhiên
----------[3] Phép cộng, phép trừ các số tự nhiên
----------[4] Phép nhân, phép chia các số tự nhiên
----------[5] Phép tính lũy thừa với số mũ tự nhiên
----------[6] Thứ tự thực hiện các phép tính
----------[7] Quan hệ chia hết. Tính chất chia hết
----------[8] Dấu hiệu chia hết cho 2, cho 5
----------[9] Dấu hiệu chia hết cho 3, cho 9
----------[0] Số nguyên tố. Hợp số
----------[A] Phân tích một số ra thừa số nguyên tố
----------[B] Ước chung và ước chung lớn nhất
----------[C] Bội chung và bội chung nhỏ nhất
-------[2] Số nguyên
----------[1] Số nguyên âm
----------[2] Tập hợp các số nguyên
----------[3] Phép cộng các số nguyên
----------[4] Phép trừ số nguyên. Quy tắc dấu ngoặc
----------[5] Phép nhân các số nguyên
----------[6] Phép chia hết hai số nguyên. Quan hệ chia hết trong tập hợp số nguyên
-------[3] Hình học trực quan
----------[1] Tam giác đều. Hình vuông. Lục giác đều
----------[2] Hình chữ nhật. Hình thoi
----------[3] Hình bình hành
----------[4] Hình thang cân
----------[5] Hình có trục đối xứng
----------[6] Hình có tâm đối xứng
----------[7] Đối xứng trong thực tiễn
-------[4] Một số yếu tố thống kê và xác suất
----------[1] Thu thập, tổ chức, biểu diễn, phân tích và xử lí dữ liệu
----------[2] Biểu đồ cột kép
----------[3] Mô hình xác suất trong một số trò chơi và thí nghiệm đơn giản
----------[4] Xác suất thực nghiệm trong một trò chơi và thí nghiệm đơn giản
-------[5] Phân số và số thập phân
----------[1] Phân số với tử và mẫu là số nguyên
----------[2] So sánh các phân số. Hỗn số dương
----------[3] Phép cộng. Phép trừ phân số
----------[4] Phép nhân, phép chia phân số
----------[5] Số thập phân
----------[6] Phép cộng, phép trừ số thập phân
----------[7] Phép nhân, phép chia số thập phân
----------[8] Ước lượng và làm tròn số
----------[9] Tỉ số. Tỉ số phần trăm
----------[0] Hai bài toán về phân số
-------[6] Hình học phẳng
----------[1] Điểm. Đường thẳng
----------[2] Hai đường thẳng cắt nhau. Hai đường thẳng song song
----------[3] Đoạn thẳng
----------[4] Tia
----------[5] Góc
----[T] Chân Trời Sáng Tạo
-------[1] Số tự nhiên
----------[1] Tập hợp. Phần tử của tập hợp
----------[2] Tập hợp số tự nhiên. Ghi số tự nhiên
----------[3] Các phép tính trong tập hợp số tự nhiên
----------[4] Lũy thừa với số mũ tự nhiên
----------[5] Thứ tự thực hiện các phép tính
----------[6] Chia hết và chia có dư. Tính chất chia hết của một tổng
----------[7] Dấu hiệu chia hết cho 2, cho 5
----------[8] Dấu hiệu chia hết cho 3, cho 9
----------[9] Ước và bội
----------[0] Số nguyên tố. Hợp số. Phân tích một số ra thừa số nguyên tố
----------[A] Hoạt động thực hành và trải nghiệm
----------[B] Ước chung. Ước chung lớn nhất
----------[C] Bội chung. Bội chung nhỏ nhất
----------[D] Hoạt động thực hành và trải nghiệm
-------[2] Số nguyên
----------[1] Số nguyên âm và tập hợp các số nguyên
----------[2] Thứ tự trong tập hợp số nguyên
----------[3] Phép cộng và phép trừ hai số nguyên
----------[4] Phép nhân và phép chia hết hai số nguyên
----------[5] Hoạt động thực hành và trải nghiệm
-------[3] Hình học trực quan và hình phẳng trong thực tiễn
----------[1] Hình vuông - Tam giác đều - Lục giác đều
----------[2] Hình chữ nhật - Hình thoi - Hình bình hành - Hình thang cân
----------[3] Chu vi và diện tích của một số hình trong thực tiễn
----------[4] Hoạt động thực hành và trải nghiệm
-------[4] Một số yếu tố thống kê
----------[1] Thu thập và phân loại dữ liệu
----------[2] Biểu diễn dữ liệu trên bảng
----------[3] Biểu đồ tranh
----------[4] Biểu đồ cột - Biểu đồ cột kép
----------[5] Hoạt động thực hành và trải nghiệm
-------[5] Phân số
----------[1] Phân số với tử số và mẫu số là số nguyên
----------[2] Tính chất cơ bản của phân số
----------[3] So sánh phân số
----------[4] Phép cộng và phép trừ phân số
----------[5] Phép nhân và phép chia phân số
----------[6] Giá trị phân số của một số
----------[7] Hỗn số
----------[8] Hoạt động thực hành và trải nghiệm
-------[6] Số thập phân
----------[1] Số thập phân
----------[2] Các phép tính với số thập phân
----------[3] Làm tròn số thập phân và ước lượng kết quả
----------[4] Tỉ số và tỉ số phần trăm
----------[5] Bài toán về tỉ số phần trăm
----------[6] Hoạt động thực hành và trải nghiệm
-------[7] Hình học trực quan
----------[1] Hình có trục đối xứng
----------[2] Hình có tâm đối xứng
----------[3] Vai trò của tính đối xứng trong thế giới tự nhiên
----------[4] Hoạt động thực hành và trải nghiệm
-------[8] Hình học phẳng và các hình hình học cơ bản
----------[1] Điểm. Đường thẳng
----------[2] Ba điểm thẳng hàng. Ba điểm không thẳng
----------[3] Hai đường thẳng cắt nhau, song song. Tia
----------[4] Đoạn thẳng. Độ dài đoạn thẳng
----------[5] Trung điểm của đoạn thẳng
----------[6] Góc
----------[7] Số đo góc. Các góc đặc biệt
----------[8] Hoạt động thực hành và trải nghiệm
-------[9] Một số yếu tố xác suất
----------[1] Phép thử nghiệm. Sự kiện
----------[2] Xác suất thực nghiệm
----------[3] Hoạt động thực hành và trải nghiệm
----[K] Kết Nối Tri Thức
-------[1] Tập hợp các số tự nhiên
----------[1] Tập hợp
----------[2] Cách ghi số tự nhiên
----------[3] Thứ tự trong tập hợp các số tự nhiên
----------[4] Phép cộng và phép trừ số tự nhiên
----------[5] Phép nhân và phép chia số tự nhiên
----------[6] Lũy thừa với số mũ tự nhiên
----------[7] Thứ tự thực hiện các phép tính
-------[2] Tính chia hết trong tập hợp các số tự nhiên
----------[8] Quan hệ chia hết và tính chất
----------[9] Dấu hiệu chia hết
----------[0] Số nguyên tố
----------[A] Ước chung. Ước chung lớn nhất
----------[B] Bội chung. Bội chung nhỏ nhất
-------[3] Số nguyên
----------[C] Tập hợp các số nguyên
----------[D] Phép cộng và phép trừ số nguyên
----------[E] Quy tắc dấu ngoặc
----------[F] Phép nhân số nguyên
----------[G] Phép chia hết. Ước và bội của một số nguyên
-------[4] Một số hình phẳng trong thực tiễn
----------[H] Hình tam giác đều. Hình vuông. Hình lục giác đều
----------[I] Hình chữ nhật. Hình thoi hình bình hành. Hình thang cân
----------[J] Chu vi và diện tích của một số tứ giác đã học
-------[5] Tính đối xứng của hình phẳng trong tự nhiên
----------[K] Hình có trục đối xứng
----------[L] Hình có tâm đối xứng
-------[6] Phân số
----------[M] Mở rộng phân số. Phân số bằng nhau
----------[N] So sánh phân số. Hỗn số dương
----------[O] Phép cộng và phép trừ phân số
----------[P] Phép nhân và phép chia phân số
----------[Q] Hai bài toán về phân số
-------[7] Số thập phân
----------[R] Số thập phân
----------[S] Tính toán với số thập phân
----------[T] Làm tròn và ước lượng
----------[U] Một số bài toán về tỉ số và tỉ số phần trăm
-------[8] Những hình học cơ bản
----------[V] Điểm và đường thẳng
----------[W] Điểm nằm giữa hai điểm. Tia
----------[X] Đoạn thẳng. Độ dài đoạn thẳng
----------[Y] Trung điểm của đoạn thẳng
----------[Z] Góc
----------[1] Số đo góc
-------[9] Dữ liệu và xác suất thực nghiệm
----------[1] Dữ liệu và thu thập dữ liệu
----------[2] Bảng thống kê và biểu đồ tranh
----------[3] Biểu đồ cột
----------[4] Biểu đồ cột kép
----------[5] Kết quả có thể và sự kiện trong trò chơi, thí nghiệm
----------[6] Xác suất thực nghiệm
----[D] Số học
-------[1] Ôn tập và bổ túc về số tự nhiên
----------[1] Tập hợp. Phần tử của tập hợp
----------[2] Tập hợp các số tự nhiên
----------[3] Ghi số tự nhiên
----------[4] Số phần tử của một tập hợp. Tập hợp con
----------[5] Phép cộng và phép nhân
----------[6] Phép trừ và phép chia
----------[7] Lũy thừa với số mũ tự nhiên. Nhân hai lũy thừa cùng cơ số.
----------[8] Chia hai lũy thừa cùng cơ số
----------[9] Thứ tự thực hiện các phép tính
----------[0] Tính chất chia hết của một tổng
----------[A] Dấu hiệu chia hết cho 2, cho 5
----------[B] Dấu hiệu chia hết cho 3, cho 9
----------[C] Ước và bội
----------[D] Số nguyên tố. Hợp số. Bảng số nguyên tố
----------[E] Phân tích một số ra thừa số nguyên tố
----------[F] Ước chung và bội chung
----------[G] Ước chung lớn nhất
----------[H] Bội chung nhỏ nhất
-------[2] Số nguyên
----------[1] Làm quen với số nguyên âm
----------[2] Tập hợp các số nguyên
----------[3] Thứ tự trong tập hợp các số nguyên
----------[4] Cộng hai số nguyên cùng dấu
----------[5] Cộng hai số nguyên khác dấu
----------[6] Tính chất của phép cộng các số nguyên
----------[7] Phép trừ hai số nguyên
----------[8] Quy tắc dấu ngoặc
----------[9] Quy tắc chuyển vế
----------[0] Nhân hai số nguyên khác dấu
----------[A] Nhân hai số nguyên cùng dấu
----------[B] Tính chất của phép nhân
----------[C] Bội và ước của một số nguyên
-------[3] Phân số
----------[1] Mở rộng khái niệm về phân số
----------[2] Phân số bằng nhau
----------[3] Tính chất cơ bản của phân số
----------[4] Rút gọn phân số
----------[5] Quy đồng mẫu số nhiều phân số
----------[6] So sánh phân số
----------[7] Phép cộng phân số
----------[8] Tính chất cơ bản của phép cộng phân số
----------[9] Phép trừ phân số
----------[0] Phép nhân phân số
----------[A] Tính chất cơ bản của phép nhân phân số
----------[B] Phép chia phân số
----------[C] Hỗn số. Số thập phân. Phần trăm
----------[D] Tìm giá trị phân số của một số cho trước
----------[E] Tìm một số biết giá trị một phân số của nó
----------[F] Tìm tỉ số của hai số
----------[G] Biểu đồ phần trăm
----[H] Hình học
-------[1] Đoạn thẳng
----------[1] Điểm. Đường thẳng
----------[2] Ba điểm thẳng hàng
----------[3] Đường thẳng đi qua hai điểm
----------[4] Thực hành: trồng cây thẳng hàng
----------[5] Tia
----------[6] Đoạn thẳng
----------[7] Độ dài đoạn thẳng
----------[8] Khi nào thì AM + MB = AB?
----------[9] Vẽ đoạn thẳng cho biết độ dài
----------[0] Trung điểm của đoạn thẳng
-------[2] Góc
----------[1] Nửa mặt phẳng
----------[2] Góc
----------[3] Số đo góc
----------[4] Khi nào góc xOy + góc yOz = góc xOz?
----------[5] Vẽ góc cho biết số đo
----------[6] Tia phân giác của góc
----------[7] Thực hành đo góc trên mặt đất
----------[8] Đường tròn
----------[9] Tam giác
-[7] Lớp 7
----[C] Cánh Diều
-------[1] Số hữu tỉ
----------[1] Tập hợp Q các số hữu tỉ
----------[2] Cộng, trừ, nhân, chia số hữu tỉ
----------[3] Phép tính luỹ thừa với số mũ tự nhiên của một số hữu tỉ
----------[4] Thứ tự thực hiện các phép tính. Quy tắc dấu ngoặc
----------[5] Biễu diễn thập phân của số hữu tỉ
-------[2] Số thực
----------[1] Số vô tỉ. Căn bậc hai số học
----------[2] Tập hợp R các số thực
----------[3] Giá trị tuyệt đối của một số thực
----------[4] Làm tròn và ước lượng
----------[5] Tỉ lệ thức
----------[6] Dãy tỉ số bằng nhau
----------[7] Đại lượng tỉ lệ thuận
----------[8] Đại lượng tỉ lệ nghịch
-------[3] Hình học trực quan
----------[1] Hình hộp chữ nhật. Hình lập phương
----------[2] Hình lăng trụ đứng tam giác. Hình lăng trụ đứng tứ giác
-------[4] Góc. Đường thẳng song song
----------[1] Góc ở vị trí đặc biệt
----------[2] Tia phân giác của một góc
----------[3] Hai đường thẳng song song
----------[4] Định lý
-------[5] Một số yếu tố thống kê và xác suất
----------[1] Thu thập, phân loại và biểu diễn dữ liệu
----------[2] Phân tích và xử lí dữ liệu
----------[3] Biểu đồ đoạn thẳng
----------[4] Biểu đồ hình quạt tròn
----------[5] Biến cố trong một số trò chơi đơn giản
----------[6] Xác suất của biến cố ngẫu nhiên trong một số trò chơi đơn giản
-------[6] Biểu thức đại số
----------[1] Biểu thức số. Biểu thức đại số
----------[2] Đa thức một biến. Nghiệm của đa thức một biến
----------[3] Phép cộng, phép trừ đa thức một biến
----------[4] Phép nhân đa thức một biến
----------[5] Phép chia đa thức một biến
-------[7] Tam giác
----------[1] Tổng các góc của một tam giác
----------[2] Quan hệ giữa góc và cạnh đối diện. Bất đẳng thức tam giác
----------[3] Hai tam giác bằng nhau
----------[4] Trường hợp bằng nhau thứ nhất của tam giác: cạnh - cạnh - cạnh
----------[5] Trường hợp bằng nhau thứ hai của tam giác: cạnh - góc - cạnh
----------[6] Trường hợp bằng nhau thứ ba của tam giác: góc - cạnh - góc
----------[7] Tam giác cân
----------[8] Đường vuông góc và đường xiên
----------[9] Đường trung trực của một đoạn thẳng
----------[0] Tính chất ba đường trung tuyến của tam giác
----------[A] Tính chất ba đường phân giác của tam giác
----------[B] Tính chất ba đường trung trực của tam giác
----------[C] Tính chất ba đường cao của tam giác
----[T] Chân Trời Sáng Tạo
-------[1] Số hữu tỉ
----------[1] Tập hợp các số hữu tỉ
----------[2] Các phép tính với số hữu tỉ
----------[3] Lũy thừa của một số hữu tỉ
----------[4] Quy tắc dấu ngoặc và quy tắc chuyển vế
----------[5] Hoạt động thực hành và trải nghiệm: Thực hành tính tiền điện
-------[2] Số thực
----------[1] Số vô tỉ. Căn bậc hai số học
----------[2] Số thực. Giá trị tuyệt đối của một số thực
----------[3] Làm tròn số và ước lượng kết quả
----------[4] Hoạt động thực hành và trải nghiệm: Tính chỉ số đánh giá thể trạng BMI (Body mass index)
-------[3] Các hình khối trong thực tiễn
----------[1] Hình hộp chữ nhật - Hình lập phương
----------[2] Diện tích xung quanh và thể tích của hình hộp chữ nhật, hình lập phương
----------[3] Hình lăng trụ đứng tam giác - Hình lăng trụ đứng tứ giác
----------[4] Diện tích xung quanh và thể tích của hình lăng trụ đứng tam giác, lăng trụ đứng tứ giác
----------[5] Hoạt động thực hành và trải nghiệm: Các bài toán về đo đạc và gấp hình
-------[4] Góc và đường thẳng song song
----------[1] Các góc ở vị trí đặc biệt
----------[2] Tia phân giác
----------[3] Hai đường thẳng song song
----------[4] Định lí và chứng minh một định lí
----------[5] Hoạt động thực hành và trải nghiệm: Vẽ hai đường thẳng song song và đo góc bằng phần mềm GeoGebra
-------[5] Một số yếu tố thống kê
----------[1] Thu thập và phân loại dữ liệu
----------[2] Biểu đồ hình quạt tròn
----------[3] Biểu đồ đoạn thẳng
----------[4] Hoạt động thực hành và trải nghiệm: Dùng biểu đồ để phân tích kết quả học tập môn Toán của lớp
-------[6] Các đại lượng tỉ lệ
----------[1] Tỉ lệ thức - Dãy tỉ số bằng nhau
----------[2] Đại lượng tỉ lệ thuận
----------[3] Đại lượng tỉ lệ nghịch
----------[4] Hoạt động thực hành và trải nghiệm: Các đại lượng tỉ lệ trong thực tế
-------[7] Biểu thức đại số
----------[1] Biểu thức số, biểu thức đại số
----------[2] Đa thức một biến
----------[3] Phép cộng và phép trừ đa thức một biến
----------[4] Phép nhân và phép chia đa thức một biến
----------[5] Hoạt động thực hành và trải nghiệm: Cách tính điểm trung bình môn học kì
-------[8] Tam giác
----------[1] Góc và cạnh của một tam giác
----------[2] Tam giác bằng nhau
----------[3] Tam giác cân
----------[4] Đường vuông góc và đường xiên
----------[5] Đường trung trực của một đoạn thẳng
----------[6] Tính chất ba đường trung trực của tam giác
----------[7] Tính chất ba đường trung tuyến của tam giác
----------[8] Tính chất ba đường cao của tam giác
----------[9] Tính chất ba đường phân giác của tam giác
----------[0] Hoạt động thực hành và trải nghiệm: Làm giàn hoa tam giác để trang trí lớp học
-------[9] Một số yếu tố xác suất
----------[1] Làm quen với biến cố ngẫu nhiên
----------[2] Làm quen với xác suất của biến cố ngẫu nhiên
----[K] Kết Nối Tri Thức
-------[1] Số hữu tỉ
----------[1] Tập hợp các số hữu tỉ
----------[2] Cộng, trừ, nhân, chia số hữu tỉ
----------[3] Lũy thừa với số mũ tự nhiên của một số hữu tỉ
----------[4] Thứ tự thực hiện các phép tính. Quy tắc chuyển vế
-------[2] Số thực
----------[5] Làm quen với số thập phân vô hạn tuần hoàn
----------[6] Số vô tỉ. Căn bậc hai số học
----------[7] Tập hợp các số thực
-------[3] Góc và đường thẳng song song
----------[8] Góc ở vị trí đặc biệt. Tia phân giác của một góc
----------[9] Hai đường thẳng song song và dấu hiệu nhận biết
----------[0] Tiên đề Euclid. Tính chất của hai đường thẳng song song
----------[A] Định lí và chứng minh định lí
-------[4] Tam giác bằng nhau
----------[B] Tổng các góc trong một tam giác
----------[C] Hai tam giác bằng nhau. Trường hợp bằng nhau thứ nhất của tam giác
----------[D] Trường hợp bằng nhau thứ hai và thứ ba của tam giác
----------[E] Các trường hợp bằng nhau của tam giác vuông
----------[F] Tam giác cân. Đường trung trực của đoạn thẳng
-------[5] Thu thập và biểu diễn dữ liệu
----------[G] Thu thập và phân loại dữ liệu
----------[H] Biểu đồ hình quạt tròn
----------[I] Biểu đồ đoạn thẳng
-------[6] Tỉ lệ thức và đại lượng tỉ lệ
----------[J] Tỉ lệ thức
----------[K] Tính chất của dãy tỉ số bằng nhau
----------[L] Đại lượng tỉ lệ thuận
----------[M] Đại lượng tỉ lệ nghịch
-------[7] Biểu thức đại số và đa thức một biến
----------[N] Biểu thức đại số
----------[O] Đa thức một biến
----------[P] Phép cộng và phép trừ đa thức một biến
----------[Q] Phép nhân đa thức một biến
----------[R] Phép chia đa thức một biến
-------[8] Làm quen với biến cố và xác suất của biến cố
----------[S] Làm quen với biến cố
----------[T] Làm quen với xác suất của biến cố
-------[9] Quan hệ giữa các yếu tố trong một tam giác
----------[U] Quan hệ giữa góc và cạnh đối diện trong một tam giác
----------[V] Quan hệ giữa đường vuông góc và đường xiên
----------[W] Quan hệ giữa ba cạnh của một tam giác
----------[X] Sự đồng quy của ba đường trung tuyến, ba đường phân giác trong một tam giác
----------[Y] Sự đồng quy của ba đường trung trực, ba đường cao trong một tam giác
-------[0] Một số hình khối trong thực tiễn
----------[1] Hình hộp chữ nhật và hình lập phương
----------[2] Hình lăng trụ đứng tam giác và hình lăng trụ đứng tứ giác
----[D] Đại số
-------[1] Số hữu tỉ. Số thực
----------[1] Tập hợp Q các số hữu tỉ
----------[2] Cộng, trừ số hữu tỉ
----------[3] Nhân, chia số hữu tỉ
----------[4] Giá trị tuyệt đối của một số hữu tỉ. Cộng, trừ, nhân, chia số thập phân
----------[5] Lũy thừa của một số hữu tỉ
----------[6] Lũy thừa của một số hữu tỉ ( tiếp theo)
----------[7] Tỉ lệ thức
----------[8] Tính chất của dãy tỉ số bằng nhau
----------[9] Số thập phân hữu hạn. Số thập phân vô hạn tuần hoàn
----------[0] Làm tròn số
----------[A] Số vô tỉ. Khái niệm về căn bậc hai
----------[B] Số thực
-------[2] Hàm số và đồ thị
----------[1] Đại lượng tỉ lệ thuận
----------[2] Một số bài toán về đại lượng tỉ lệ thuận
----------[3] Đại lượng tỉ lệ nghịch
----------[4] Một số bài toán về đại lượng tỉ lệ nghịch
----------[5] Hàm số
----------[6] Mặt phẳng toạ độ
----------[7] Đồ thị hàm số y = ax (a # 0)
-------[3] Thống kê
----------[1] Thu thập số liệu thống kê, tần số
----------[2] Bảng tần số các giá trị của dấu hiệu
----------[3] Biểu đồ
----------[4] Số trung bình cộng
-------[4] Biểu thức đại số
----------[1] Khái niệm về biểu thức đại số
----------[2] Giá trị của một biểu thức đại số
----------[3] Đơn thức
----------[4] Đơn thức đồng dạng
----------[5] Đa thức
----------[6] Cộng, trừ đa thức
----------[7] Đa thức một biến
----------[8] Cộng, trừ đa thức một biến
----------[9] Nghiệm của đa thức một biến
----[H] Hình học
-------[1] Đường thẳng vuông góc. Đường thẳng song song
----------[1] Hai góc đối đỉnh
----------[2] Hai đường thẳng vuông góc
----------[3] Các góc tạo bởi một đường thẳng cắt hai đường thẳng
----------[4] Hai đường thẳng song song
----------[5] Tiên đề Ơ-clit về đường thẳng song song
----------[6] Từ vuông góc đến song song
----------[7] Định lí
-------[2] Tam giác
----------[1] Tổng ba góc của một tam giác
----------[2] Hai tam giác bằng nhau
----------[3] Trường hợp bằng nhau thứ nhất của tam giác cạnh - cạnh - cạnh (c.c.c)
----------[4] Trường hợp bằng nhau thứ hai của tam giác cạnh - góc - cạnh (c.g.c)
----------[5] Trường hợp bằng nhau thứ ba của tam giác góc - cạnh - góc (g.c.g)
----------[6] Tam giác cân
----------[7] Định lí Py-ta-go
----------[8] Các trường hợp bằng nhau của tam giác vuông
-------[3] Quan hệ giữa các yểu tố trong tam giác. Các đường đồng quy của tam giác
----------[1] Quan hệ giữa góc và cạnh đối diện trong một tam giác
----------[2] Quan hệ giữa đường vuông góc và đường xiên, đường xiên và hình chiếu
----------[3] Quan hệ giữa ba cạnh của một tam giác, bất đẳng thức tam giác
----------[4] Tính chất ba đường trung tuyến của tam giác
----------[5] Tính chất tia phân giác của một góc
----------[6] Tính chất ba đường phân giác của tam giác
----------[7] Tính chất đường trung trực của một đoạn thẳng
----------[8] Tính chất ba đường trung trực của tam giác
----------[9] Tính chất ba đường cao của tam giác
-[8] Lớp 8
----[C] Cánh Diều
-------[1] Đa thức nhiều biến
----------[1] Đơn thức nhiều biến. Đa thức nhiều biến
----------[2] Các phép tính với đa thức nhiều biến
----------[3] Hằng đẳng thức đáng nhớ
----------[4] Vận dụng hằng đẳng thức vào phân tích đa thức thành nhân tử
-------[2] Phân thức đại số
----------[1] Phân thức đại số
----------[2] Phép cộng, phép trừ phân thức đại số
----------[3] Phép nhân, phép chia phân thức đại số
-------[3] Hàm số và đồ thị
----------[1] Hàm số
----------[2] Mặt phẳng tọa độ. Đồ thị của hàm số
----------[3] Hàm số bậc nhất y = ax + b (a khác 0)
----------[4] Đồ thị của hàm số bậc nhất y = ax + b (a khác 0)
-------[4] Hình học trực quan
----------[1] Hình chóp tam giác đều
----------[2] Hình chóp tứ giác đều
-------[5] Tam giác. Tứ giác
----------[1] Định lí Pythagore
----------[2] Tứ giác
----------[3] Hình thang cân
----------[4] Hình bình hành
----------[5] Hình chữ nhật
----------[6] Hình thoi
----------[7] Hình vuông
-------[6] Một số yếu tố thống kê và xác suất
----------[1] Thu thập và phân loại dữ liệu
----------[2] Mô tả và biểu diễn dữ liệu trên các bảng, biểu đồ
----------[3] Phân tích và xử lí dữ liệu thu được ở dạng bảng, biểu đồ
----------[4] Xác suất của biến cố ngẫu nhiên trong một số trò chơi đơn giản
----------[5] Xác suất thực nghiệm của một biến cố trong một số trò chơi đơn giản
-------[7] Phương trình bậc nhất một ẩn
----------[1] Phương trình bậc nhất một ẩn
----------[2] Ứng dụng của phương trình bậc nhất một ẩn
-------[8] Tam giác đồng dạng. Hình đồng dạng
----------[1] Định lí Thalès trong tam giác
----------[2] Ứng dụng của định lí Thalès trong tam giác
----------[3] Đường trung bình của tam giác
----------[4] Tính chất đường phân giác của tam giác
----------[5] Tam giác đồng dạng
----------[6] Trường hợp đồng dạng thứ nhất của tam giác
----------[7] Trường hợp đồng dạng thứ hai của tam giác
----------[8] Trường hợp đồng dạng thứ ba của tam giác
----------[9] Hình đồng dạng
----------[0] Hình đồng dạng trong thực tiễn
----[T] Chân Trời Sáng Tạo
-------[1] Biểu thức đại số
----------[1] Đơn thức và đa thức nhiều biến
----------[2] Các phép toán với đa thức nhiều biến
----------[3] Hằng đẳng thức đáng nhớ
----------[4] Phân tích đa thức thành nhân tử
----------[5] Phân thức đại số
----------[6] Cộng, trừ phân thức
----------[7] Nhân, chia phân thức
-------[2] Các hình khối trong thực tiễn
----------[1] Hình chóp tam giác đều – Hình chóp tứ giác đều
----------[2] Diện tích xung quanh và thể tích của hình chóp tam giác đều, hình chóp tứ giác đều
-------[3] Định lí Pythagore. Các loại tứ giác thường gặp
----------[1] Định lí Pythagore
----------[2] Tứ giác
----------[3] Hình thang – Hình thang cân
----------[4] Hình bình hành – Hình thoi
----------[5] Hình chữ nhật – Hình vuông
-------[4] Một số yếu tố thống kê
----------[1] Thu thập và phân loại dữ liệu
----------[2] Lựa chọn dạng biểu đồ để biểu diễn dữ liệu
----------[3] Phân tích dữ liệu
-------[5] Hàm số và đồ thị
----------[1] Khái niệm hàm số
----------[2] Toạ độ của một điểm và đồ thị của hàm số
----------[3] Hàm số bậc nhất y = ax + b (a khác 0)
----------[4] Hệ số góc của đường thẳng
-------[6] Phương trình
----------[1] Phương trình bậc nhất một ẩn
----------[2] Giải bài toán bằng cách lập phương trình bậc nhất
-------[7] Định lí Thalès
----------[1] Định lí Thalès trong tam giác
----------[2] Đường trung bình của tam giác
----------[3] Tính chất đường phân giác của tam giác
-------[8] Hình đồng dạng
----------[1] Hai tam giác đồng dạng
----------[2] Các trường hợp đồng dạng của hai tam giác
----------[3] Các trường hợp đồng dạng của hai tam giác vuông
----------[4] Hai hình đồng dạng
-------[9] Một số yếu tố xác suất
----------[1] Mô tả xác suất bằng tỉ số
----------[2] Xác suất lí thuyết và xác suất thực nghiệm
----[K] Kết Nối Tri Thức
-------[1] Đa thức
----------[1] Đơn thức
----------[2] Đa thức
----------[3] Phép cộng và phép trừ đa thức
----------[4] Phép nhân đa thức
----------[5] Phép chia đa thức cho đơn thức
-------[2] Hằng đẳng thức đáng nhớ và ứng dụng
----------[6] Hiệu hai bình phương. Bình phương của một tổng hay một hiệu
----------[7] Lập phương của một tổng hay một hiệu
----------[8] Tổng và hiệu hai lập hương
----------[9] Phân tích đa thức thành nhân tử
-------[3] Tứ giác
----------[0] Tứ giác
----------[A] Hình thang cân
----------[B] Hình bình hành
----------[C] Hình chữ nhật
----------[D] Hình thoi và hình vuông
-------[4] Định lí Thalès
----------[E] Định lí Thalès trong tam giác
----------[F] Đường trung bình của tam giác
----------[G] Tính chất đường phân giác của tam giác
-------[5] Dữ liệu và biểu đồ
----------[H] Thu thập và phân loại dữ liệu
----------[I] Biểu diễn dữ liệu bằng bảng, biểu đồ
----------[J] Phân tích số liệu thống kê dựa vào biểu đồ
-------[6] Phân thức đại số
----------[K] Phân thức đại số
----------[L] Tính chất cơ bản của phân thức đại số
----------[M] Phép cộng và phép trừ phân thức đại số
----------[N] Phép nhân và phép chia phân thức đại số
-------[7] Phương trình bậc nhất và hàm số bậc nhất
----------[O] Phương trình bậc nhất một ẩn
----------[P] Giải bài toán bằng cách lập phương trình
----------[Q] Khái niệm hàm số và đồ thị của hàm số
----------[R] Hàm số bậc nhất và đồ thị của hàm số bậc nhất
----------[S] Hệ số góc của đường thẳng
-------[8] Mở đầu về tính xác suất của biến cố
----------[T] Kết quả có thể và kết quả thuận lợi
----------[U] Cách tính xác suất của biến cố bằng tỉ số
----------[V] Mối liên hệ giữa xác suất thực nghiệm với xác suất và ứng dụng
-------[9] Tam giác đồng dạng
----------[W] Hai tam giác đồng dạng
----------[X] Ba trường hợp đồng dạng của hai tam giác
----------[Y] Định lí Pythagore và ứng dụng
----------[Z] Các trường hợp đồng dạng của hai tam giác vuông
----------[1] Hình đồng dạng
-------[0] Một số hình khối trong thực tiễn
----------[2] Hình chóp tam giác đều
----------[3] Hình chóp tứ giác đều
----[D] Đại số
-------[1] Phép nhân và phép chia đa thức
----------[1] Nhân đơn thức với đa thức
----------[2] Nhân đa thức với đa thức
----------[3] Những hằng đẳng thức đáng nhớ
----------[4] Những hằng đẳng thức đáng nhớ (tiếp)
----------[5] Những hằng đẳng thức đáng nhớ (tiếp)
----------[6] Phân tích đa thức thành nhân tử bằng phương pháp đặt nhân tử chung
----------[7] Phân tích đa thức thành nhân tử bằng phương pháp dùng hằng đẳng thức
----------[8] Phân tích đa thức thành nhân tử bằng phương pháp nhóm hạng tử
----------[9] Phân tích đa thức thành nhân tử bằng cách phối hợp nhiều phương pháp
----------[0] Chia đơn thức cho đơn thức
----------[A] Chia đa thức cho đơn thức
----------[B] Chia đa thức một biến đã sắp xếp
-------[2] Phân thức đại số
----------[1] Phân thức đại số
----------[2] Tính chất cơ bản của phân thức
----------[3] Rút gọn phân thức
----------[4] Quy đồng mẫu thức nhiều phân thức
----------[5] Phép cộng các phân thức đại số
----------[6] Phép trừ các phân thức đại số
----------[7] Phép nhân các phân thức đại số
----------[8] Phép chia các phân thức đại số
----------[9] Biến đổi các biểu thức hữu tỉ. Giá trị của phân thức
-------[3] Phương trình bậc nhất một ẩn
----------[1] Mở đầu về phương trình
----------[2] Phương trình bậc nhất một ẩn và cách giải
----------[3] Phương trình đưa được về dạng ax + b = 0
----------[4] Phương trình tích
----------[5] Phương trình chứa ẩn ở mẫu
----------[6] Giải bài toán bằng cách lập phương trình
----------[7] Giải bài toán bằng cách lập phương trình (tiếp)
-------[4] Bất phương trình bậc nhất một ẩn
----------[1] Liên hệ giữa thứ tự và phép cộng
----------[2] Liên hệ giữa thứ tự và phép nhân
----------[3] Bất phương trình một ẩn
----------[4] Bất phương trình bậc nhất một ẩn
----------[5] Phương trình chứa dấu giá trị tuyệt đối
----[H] Hình học
-------[1] Tứ giác
----------[1] Tứ giác
----------[2] Hình thang
----------[3] Hình thang cân
----------[4] Đường trung bình của tam giác, của hình thang
----------[5] Dựng hình bằng thước và compa. Dựng hình thang
----------[6] Đối xứng trục
----------[7] Hình bình hành
----------[8] Đối xứng tâm
----------[9] Hình chữ nhật
----------[0] Đường thẳng song song với một đường thẳng cho trước
----------[A] Hình thoi
----------[B] Hình vuông
-------[2] Đa giác, diện tích đa giác
----------[1] Đa giác. Đa giác đều
----------[2] Diện tích hình chữ nhật
----------[3] Diện tích tam giác
----------[4] Diện tích hình thang
----------[5] Diện tích hình thoi
----------[6] Diện tích đa giác]
-------[3] Tam giác đồng dạng
----------[1] Định lí Ta-lét trong tam giác
----------[2] Định lí đảo và hệ quả của định lí Ta-lét
----------[3] Tính chất đường phân giác của tam giác
----------[4] Khái niệm hai tam giác đồng dạng
----------[5] Trường hợp đồng dạng thứ nhất
----------[6] Trường hợp đồng dạng thứ hai
----------[7] Trường hợp đồng dạng thứ ba
----------[8] Các trường hợp đồng dạng của tam giác vuông
----------[9] Ứng dụng thực tế của tam giác đồng dạng
-------[4] Hình lăng trụ đứng. Hình chóp đều
----------[1] Hình hộp chữ nhật
----------[2] Hình hộp chữ nhật (tiếp)
----------[3] Thể tích của hình hộp chữ nhật
----------[4] Hình lăng trụ đứng
----------[5] Diện tích xung quanh của hình lăng trụ đứng
----------[6] Thể tích của hình lăng trụ đứng
----------[7] Hình chóp đều và hình chóp cụt đều
----------[8] Diện tích xung quanh của hình chóp
----------[9] Thể tích của hình chóp đều
-[9] Lớp 9
----[D] Đại số
-------[1] Căn bậc hai. Căn bậc ba
----------[1] Căn bậc hai
----------[2] Căn thức bậc hai và hằng đẳng thức
----------[3] Liên hệ giữa phép nhân và phép khai phương
----------[4] Liên hệ giữa phép chia và phép khai phương
----------[5] Bảng căn bậc hai
----------[6] Biến đổi đơn giản biểu thức chứa căn thức bậc hai
----------[7] Biến đổi đơn giản biểu thức chứa căn thức bậc hai (tiếp theo)
----------[8] Rút gọn biểu thức chứa căn bậc hai
----------[9] Căn bậc ba
-------[2] Hàm số bậc nhất
----------[1] Nhắc lại và bổ sung các khái niệm về hàm số
----------[2] Hàm số bậc nhất
----------[3] Đồ thị của hàm số y = ax + b (a ≠ 0)
----------[4] Đường thẳng song song và đường thẳng cắt nhau
----------[5] Hệ số góc của đường thẳng y = ax + b (a ≠ 0)
-------[3] Hệ hai phương trình bậc nhất hai ẩn
----------[1] Phương trình bậc nhất hai ẩn
----------[2] Hệ hai phương trình bậc nhất hai ẩn
----------[3] Giải hệ phương trình bằng phương pháp thế
----------[4] Giải hệ phương trình bằng phương pháp cộng đại số
----------[5] Giải bài toán bằng cách lập hệ phương trình
----------[6] Giải bài toán bằng cách lập hệ phương trình (tiếp theo)
-------[4] Hàm số y = ax^2 (a ≠ 0). Phương trình bậc hai một ẩn
----------[1] Hàm số y = ax^2 (a ≠ 0)
----------[2] Đồ thị của hàm số y = ax^2 (a ≠ 0)
----------[3] Phương trình bậc hai một ẩn
----------[4] Công thức nghiệm của phương trình bậc hai
----------[5] Công thức nghiệm thu gọn
----------[6] Hệ thức vi-ét và ứng dụng
----------[7] Phương trình quy về phương trinh bậc hai
----------[8] Giải bài toán bằng cách lập phương trình
----[H] Hình học
-------[1] Hệ thức lượng trong tam giác vuông
----------[1] Một số hệ thức về cạnh và đường cao trong tam giác vuông
----------[2] Tỉ số lượng giác của góc nhọn
----------[3] Bảng lượng giác
----------[4] Một số hệ thức về cạnh và góc trong tam giác vuông
----------[5] Ứng dụng thực tế các tỉ số lượng giác của góc nhọn. Thực hành ngoài trời
-------[2] Đường tròn
----------[1] Sự xác định của đường tròn. Tính chất đối xứng của đường tròn
----------[2] Đường kính và dây của đường tròn
----------[3] Liên hệ giữa dây và khoảng cách từ tâm đến dây
----------[4] Vị trí tương đối của đường thẳng và đường tròn
----------[5] Dấu hiệu nhận biết tiếp tuyến của đường tròn
----------[6] Tính chất của hai tiếp tuyến cắt nhau
----------[7] Vị trí tương đối của hai đường tròn
----------[8] Vị trí tương đối của hai đường tròn (tiếp theo)
-------[3] Góc với đường tròn
----------[1] Góc ở tâm. Số đo cung
----------[2] Liên hệ giữa cung và dây
----------[3] Góc nội tiếp
----------[4] Góc tạo bởi tia tiếp tuyến và dây cung
----------[5] Góc có đỉnh ở bên trong đường tròn. Góc có đỉnh ở bên ngoài đường tròn
----------[6] Cung chứa góc
----------[7] Tứ giác nội tiếp
----------[8] Đường tròn ngoại tiếp. Đường tròn nội tiếp
----------[9] Độ dài đường tròn, cung tròn
----------[0] Diện tích hình tròn, hình quạt tròn
-------[4] Hình trụ - hình nón - hình cầu
----------[1] Hình trụ - diện tích xung quanh và thể tích hình trụ
----------[2] Hình nón - hình nón cụt. Diện tích xung quanh và thể tích của hình nón, hình nón cụt
----------[3] Hình cầu. Diện tích hình cầu và thể tích hình cầu
%Kết thúc nội dung ID
